\documentclass[../main.tex]{subfiles}

\begin{document}
% -----------------------------
\section{Creating new dialect}
% -----------------------------
Add imports and define custom dialect.
\begin{lstlisting} [language=Python, caption={Define new dialect}]
from xdsl.context import Context
from xdsl.ir import Dialect
from xdsl.dialects.builtin import Builtin
from xdsl.dialects import func


# For now: Empty HiCompiler dialect (no ops / attributes yet)
HiCompiler = Dialect(
    "hc",      # dialect name: will give ops like `hc.foo` later
    (),        # operations (we will add ops here in later steps)
    (),        # attributes / types (also later)
)
\end{lstlisting}
The point of this step is just is just to see how can we register and load the dialect.
\begin{lstlisting} [language=Python, caption={Define new dialect}]
def main() -> None:
    ctx = Context()

    # Load the built-in dialects
    ctx.load_dialect(Builtin)
    ctx.load_dialect(func.Func)

    # Now load your custom HiCompiler dialect instance
    ctx.load_dialect(HiCompiler)

    # The point is just: 'can we register and load the dialect without errors?'
    print(ctx)
    print("\nHiCompiler dialect registered and loaded successfully.")

if __name__ == "__main__":
    main()
\end{lstlisting}

\end{document}
